% ============================================
% Assistant Professor Cover Letter – Loyola University Chicago (Economics)
% ============================================

\documentclass[11pt,letterpaper]{letter}

\usepackage{geometry}
\usepackage{microtype}
\usepackage[hidelinks]{hyperref}

\geometry{left=25mm,right=25mm,top=25mm,bottom=25mm}

\signature{Decory Edwards}
\address{
Department of Economics \\
Johns Hopkins University \\
Baltimore, MD 21218 \\
\texttt{dedwar65@jh.edu}
}

\begin{document}

\begin{letter}{Search Committee \\
Department of Economics \\
Quinlan School of Business \\
Loyola University Chicago \\
Chicago, IL \\ USA}

\opening{Dear Members of the Search Committee,}

I am writing to apply for the tenure-track Assistant Professor position in Economics in the Quinlan School of Business at Loyola University Chicago. I am a Ph.D.\ candidate in Economics at Johns Hopkins University, advised by Professors Christopher D.~Carroll, Nicholas W.~Papageorge, and Stelios Fourakis, and I expect to receive my degree in May~2026. My research fields are macroeconomics, wealth and income dynamics, and computational economics.

My research agenda focuses on understanding the determinants of wealth inequality and the mechanisms through which heterogeneity in returns to wealth shapes aggregate outcomes. In my job-market paper, I estimate the degree of heterogeneity in individual portfolio returns using micro-data from the Survey of Consumer Finances and simulate a structural model in which households face idiosyncratic return risk. I show that allowing for persistent heterogeneity in returns substantially improves the model’s ability to match the empirical distribution of wealth, offering a tractable way to connect micro-level return dispersion to macroeconomic inequality.

In a second project, I use the Health and Retirement Study to examine the relationship between interpersonal trust and portfolio performance. Building on the literature that finds a hump-shaped relationship between trust and income, I show that a similar relationship holds for individual returns to wealth measured in the HRS data, highlighting a behavioral channel through which social attitudes shape saving and investment outcomes. This work also provides new estimates of the persistent component of returns to net worth, which motivate the calibration of heterogeneity in my structural model.

\textbf{Teaching.} I have experience teaching and supporting courses in \textit{Principles of Macroeconomics}, \textit{Intermediate Macroeconomic Theory}, and upper-level macroeconomic policy topics. My teaching emphasizes clarity, economic intuition, and the disciplined use of data to connect theory to applied questions in finance and business. At Loyola, I would be well prepared to teach courses in \textit{Money and Banking} and \textit{Monetary Economics}, as well as introductory and intermediate macroeconomics for business students. I am particularly interested in helping students understand how monetary policy, financial institutions, and asset markets interact with real economic outcomes.

Loyola University Chicago’s Quinlan School of Business is an especially strong environment for my work. The School’s emphasis on rigorous economics training within a business curriculum, its engagement with real-world financial and policy questions, and its Jesuit commitment to ethical leadership align well with my approach to research and teaching. I am drawn to Loyola’s focus on educating students to think critically about economic systems and to apply analytical tools responsibly in professional settings. I would welcome the opportunity to contribute to the economics programs at both the undergraduate and graduate levels and to participate actively in the intellectual and service-oriented life of the School.

Thank you very much for your time and consideration. I would welcome the opportunity to discuss my application further.

\closing{Sincerely,}

\end{letter}
\end{document}
